% The Slide Definitions
%document
\documentclass[10pt]{beamer}
%theme
\usetheme{metropolis}
% packages
\usepackage{color}
\usepackage{listings}
\usepackage[ngerman]{babel}
\usepackage[utf8]{inputenc}
\usepackage{multicol}


% color definitions
\definecolor{mygreen}{rgb}{0,0.6,0}
\definecolor{mygray}{rgb}{0.5,0.5,0.5}
\definecolor{mymauve}{rgb}{0.58,0,0.82}

\lstset{
    backgroundcolor=\color{white},
    % choose the background color;
    % you must add \usepackage{color} or \usepackage{xcolor}
    basicstyle=\footnotesize\ttfamily,
    % the size of the fonts that are used for the code
    breakatwhitespace=false,
    % sets if automatic breaks should only happen at whitespace
    breaklines=true,                 % sets automatic line breaking
    captionpos=b,                    % sets the caption-position to bottom
    commentstyle=\color{mygreen},    % comment style
    % deletekeywords={...},
    % if you want to delete keywords from the given language
    extendedchars=true,
    % lets you use non-ASCII characters;
    % for 8-bits encodings only, does not work with UTF-8
    frame=single,                    % adds a frame around the code
    keepspaces=true,
    % keeps spaces in text,
    % useful for keeping indentation of code
    % (possibly needs columns=flexible)
    keywordstyle=\color{blue},       % keyword style
    % morekeywords={*,...},
    % if you want to add more keywords to the set
    numbers=left,
    % where to put the line-numbers; possible values are (none, left, right)
    numbersep=5pt,
    % how far the line-numbers are from the code
    numberstyle=\tiny\color{mygray},
    % the style that is used for the line-numbers
    rulecolor=\color{black},
    % if not set, the frame-color may be changed on line-breaks
    % within not-black text (e.g. comments (green here))
    stepnumber=1,
    % the step between two line-numbers.
    % If it's 1, each line will be numbered
    stringstyle=\color{mymauve},     % string literal style
    tabsize=4,                       % sets default tabsize to 4 spaces
    % show the filename of files included with \lstinputlisting;
    % also try caption instead of title
    language = Python,
	showspaces = false,
	showtabs = false,
	showstringspaces = false,
	escapechar = ,
}

\def\ContinueLineNumber{\lstset{firstnumber=last}}
\def\StartLineAt#1{\lstset{firstnumber=#1}}
\let\numberLineAt\StartLineAt



\newcommand{\codeline}[1]{
	\alert{\texttt{#1}}
}


% Author and Course information
% This Document contains the information about this course.

% Authors of the slides
\author{Felix Döring, Felix Wittwer, Anton Obersteiner}

% Name of the Course
\institute{Python-Kurs}

% Fancy Logo 
\titlegraphic{\hfill\includegraphics[height=1.25cm]{../templates/fsr_logo_cropped}}



% Custom Bindings
% \newcommand{\codeline}[1]{
%	\alert{\texttt{#1}}
%}


% Presentation title
\title{Module, Packete und PIP}
\date{\today}


\begin{document}

\maketitle

\begin{frame}{Gliederung}
	\begin{multicols}{2}
		\setbeamertemplate{section in toc}[sections numbered]
		\tableofcontents
	\end{multicols}
\end{frame}

% ---------------------------------- Module ----------------------------------

\section{Module}
\subsection{Eigene Module}
\begin{frame}[fragile]{Module}
	\begin{itemize}
		\item Ein Modul ist die python-interne Repr\"asentation einer \texttt{.py} Datei
		\item Der Dateiname setzt sich daher aus Modulname + \texttt{.py} zusammen
		\item Module beinhalten eine beliebige Anzahl an Definitionen (Klassen, Funktionen, Variablen/Konstanten)
		\item Das komplette Modul wird mit \alert{\texttt{import}} hinzugef\"ugt
		\item Einzelne Inhalte mit \alert{\texttt{from}} Modul \alert{\texttt{import}} Name
		\item Mit Hilfe von \alert{\texttt{as}} kann ein Alias f\"ur den importierten Namen angelegt werden
	\end{itemize}
\end{frame}

\begin{frame}{Ein Beispiel}
	Die Datei \textbf{incdec.py} enth\"alt eine Reihe von Funktionen:
	\lstinputlisting{resources/04_modules_and_packages/incdec.py}
\end{frame}

\begin{frame}{Ein Beispiel}
	In einem anderen Python-Script kann ich diese Funktionen nutzen:
	\lstinputlisting[lastline=15]{resources/04_modules_and_packages/modules.py}
\end{frame}

\begin{frame}{Ein Beispiel}
	In einem anderen Python-Script kann ich diese Funktionen nutzen:
	\lstinputlisting[firstline=17]{resources/04_modules_and_packages/modules.py}
\end{frame}

\subsection{Der Sinn der Boilerplate}
\begin{frame}{Der Sinn der Boilerplate}
	\begin{itemize}
		\item Verhindert, dass Code beim Importieren eines Scriptes ausgef\"uhrt wird
		\item Beim Aufruf des Moduls \"uber \texttt{import} ist der Name des Moduls \alert{nicht \texttt{\_\_main\_\_}}, sondern der Name des Scriptes (im vorherigen Beispiel w\"are das \texttt{incdec})
	\end{itemize}
\end{frame}



\subsection{Suchpfad f\"ur Module}
\begin{frame}{Suchpfad f\"ur Module}
	\begin{itemize}
		\item Python sucht Module beim Import an allen im \alert{\textbf{PYTHONPATH}} aufgelisteten Ordnern
		\item dieser findet sich in Python unter \alert{\texttt{sys.path}}, in der Kommandozeile in der Umgebungsvariable PYTHONPATH
		\item \texttt{sys.path} l\"asst sich zur Laufzeit im Interpreter \"andern (z.B. um nachtr\"aglich neue eigene Module hinzuzuf\"ugen)
	\end{itemize}
\end{frame}

\begin{frame}{Suchpfad f\"ur Module}
	Standardm\"a\ss{}ig enth\"alt der Suchpfad folgende Module: \\ \ \\
	\begin{itemize}
		\item die \textit{Standardbibliothek} der derzeitig verwendeten Python-Version
		\item das \textit{aktuelle Verzeichnis}, in dem der Interpreter aufgerufen wurde
		\item eine \textit{Verzeichnis mit plattfomspezifischen Modulen}, z.B.
			"/usr/local/Cellar/python3/3.5.1/Frameworks/Python.framework/
			Versions/3.5/lib/python3.5/plat-darwin" f\"ur Mac
		\item \textit{Benutzerspezifische Module f\"ur die jeweilige Python Version}, z.B. "/usr/local/lib/python3.5/site-packages" f\"ur Python 3.5
	\end{itemize}
\end{frame}



\subsection{Suchpfad modifizieren}
\begin{frame}{Suchpfad modifizieren}
	Zus\"atzliche Verzeichnisse f\"ur den PYTHONPATH kann man beim Aufruf \"ubergeben:\\
	\hspace*{.5cm}\codeline{\$ PYTHONPATH=/my/directory python3 script.py}\\ \ \\
	Oder im Programm:
	\lstinputlisting{resources/04_modules_and_packages/path.py}
	\"Anderungen am Pfad sind erst \textbf{nach} der entsprechenden Codezeile verf\"ugbar.
\end{frame}



\subsection{Standardmodule}
\begin{frame}{Standardmodule}
	Python liefert viele n\"utzliche Module bereits in der Standardbibliothek mit, z.B.:
	\texttt{sys}, \texttt{os}, \texttt{http}, \texttt{re}, \texttt{fuctools}, \texttt{itertools}, \texttt{collections}, \texttt{hashlib}, \texttt{urllib} und viele mehr. \\ \ \\
	Daher bezeichnet man Python oft auch als \textbf{Batteries included}.
\end{frame}

% --------------------------------- Packages ---------------------------------

\section{Packages}
\begin{frame}{Packages}
	\begin{itemize}
		\item Packages sind Ordner, die mindestens ein \alert{\texttt{\_\_init\_\_.py}} Modul enthalten.
		\item der Inhalt dieses Moduls ist prinzipiell egal
		\item Packages k\"onnen genau wie Module importiert werden
		\item Wird das Package selbst importiert, sind alle Definitionen in \texttt{\_\_init\_\_.py} \"uber das importierte Package erreichbar.
		\item wird benutzt, um Module in sinnvolle Gruppen zusammenzufassen
	\end{itemize}
\end{frame}

% -------------------- merged from presentation about 'pip' -------------------

\section{PIP}
\subsection{Installation von PIP}
\begin{frame}{Installation von PIP}
	Nicht jedes Modul ist in der Standardbibliothek vorhanden, muss aber nicht zwingen manuell als Datei hinzugef\"ugt werden. Man kann ganz einfach PIP zur Hilfe nehmen.\\[.5cm]
	\textbf{Installation:}\\
	\begin{enumerate}
		\item Downloaden der \texttt{get-pip.py} von pypa \href{https://bootstrap.pypa.io/get-pip.py}{(hier)}
		\item Installieren via \texttt{python3 get-pip.py}\\[.5cm]
	\end{enumerate}

	Hinweis f\"ur Mac-Nutzer:\\
	Bei der Pythoninstallation via \textit{Homebrew} wird pip gleich mit installiert.
\end{frame}



\subsection{Verwendung}
\begin{frame}{Verwendung}
	Die Installation von Modulen ist einfach und es gibt mehrere M\"oglichkeiten:
	\begin{itemize}
		\item die aktuellste Version eines Moduls:\\
		\noindent\hspace*{.5cm} \texttt{pip install Modulname}
		\item eine bestimmte Version:\\
		\noindent\hspace*{.5cm} \texttt{pip install Modulname == 2.1}
		\item alle Versionen ab einem Minimum:\\
		\noindent\hspace*{.5cm} \texttt{pip install Modulname >= 2.1}
		\item alle Versionen bis zu einem Maximum:\\
		\noindent\hspace*{.5cm} \texttt{pip install Modulname <= 2.1}
	\end{itemize}
\end{frame}



\subsection{Die requirements.txt}
\begin{frame}{Die requirements.txt}
	Wenn jemand Scripte ausf\"uhren m\"ochte, in denen Module verwendet werden, die mit Hilfe von PIP installiert wurden, ist es hilfreich eine \alert{\texttt{requirements.txt}} mitzuliefern.\\
	In dieser sind alle Module aufgelistet, die ben\"otigt werden.\\[.5cm]
	\lstinputlisting{resources/05_pip/requirements.txt}

	Via \alert{\texttt{pip install -r requirements.txt}} k\"onnen alle Module von der Liste installiert werden
\end{frame}

\end{document}
